\documentclass{resume} % Use the custom resume.cls style
\usepackage{times}
\usepackage{paralist}

\usepackage[left=0.5in,top=0.2in,right=0.5in,bottom=0.5in]{geometry} % Adjusted margins
\newcommand{\tab}[1]{\hspace{.2667\textwidth}\rlap{#1}} 
\newcommand{\itab}[1]{\hspace{0em}\rlap{#1}}
\name{Liam Cawley} % Your name
\address{+1 (917) 319-9518 \\ New York, NY \\ \href{mailto:cawleyl@umich.edu}{cawleyl@umich.edu} \\ github.com/fxcawley}
\usepackage{setspace}

\begin{document}
\begin{spacing}{0.9}

%----------------------------------------------------------------------------------------
%	EDUCATION SECTION
%----------------------------------------------------------------------------------------
\vspace{-0.5em}


\begin{rSection}{Education}
{\bf UNIVERSITY OF MICHIGAN}, {\bf College of Engineering} \hfill {May 2025}\\
{\bf Bachelor of Engineering}, Computer Science Engineering, Minor in Mathematics
\end{rSection}
\vspace{-0.75em}

%----------------------------------------------------------------------------------------
%	RELEVANT COURSEWORK
%----------------------------------------------------------------------------------------

\begin{rSection}{Coursework}
\vspace{-0.25em}
\begin{inparaitem}
    \item Random Numerical Linear Algebra
    \item Principles of Machine Learning
    \item Science of Language Models
    \item Game Theory
    \item Linear Algebra
\end{inparaitem}
\vspace{-0.5em}

\begin{inparaitem}
    \item Graph Algorithms
    \item Mathematics of Finance
    \item Probability Theory
    \item Randomness and Computation
    \item Matrix Algebra
    \item Computer Vision 
\end{inparaitem}
\end{rSection}
\vspace{-0.75em}


\begin{rSection}{Experience}
\vspace{-0.5em}


\textbf{Machine Learning Engineer} \hfill May 2025 -- Present \\
KLA Corporation \hfill \textit{Ann Arbor, MI}
\vspace{-0.5em}

\begin{itemize}\itemsep 0pt
    \item Led backend implementation of a production ETL and inference pipeline (DAL) that served proprietary encrypted vision models as a child of a legacy 12-microservice, 8 compute node Docker swarm (DART) deployed to on prem 12-GPU rack. 
    
    \item Developed a hybrid local/remote inter-process communication model, utilizing POSIX Shared Memory (/dev/shm) for zero-copy local fast-path IPC on single-node GPU workers, and gRPC for cross-node inter-service communication.

    \item Led release engineering for target-node deployments, including appsettings configuration, dependency validation, service health checks, and CI/CD orchestration across Docker Swarm nodes.
    
    \item Built an AI analysis service with Gemma-4-26B-A4B served via  vLLM container with TensorRT (orchestrated with LangGraph) with multiple subagents, RAG, in-memory session management, and 20 integrated tools with realtime DART service integration. 
    
    \item Awarded 1st Place at KLA's global algorithms conference for developing a Physics-Informed Neural Network system and presenting its technical roadmap to senior division leadership.
        

\end{itemize}
\vspace{-0.25em}

\textbf{High Performance Computing Intern} \hfill May 2024 -- Aug 2024 \\
KLA Corporation \hfill \textit{San Jose, CA}
\vspace{-0.5em}

\begin{itemize}\itemsep 0pt
  \item Built a proprietary vision model system as a proof of the hypothesis that physical chip layout geometries could be associated with systematic occurrences of atomic-scale transistor defects.
 
  \item Co-authored a high-performance design-geometry parser by extending KLA’s internal \textbf{GDSTK} C++ engine, optimizing the extraction and localization of layout structures for specific wafer defect locations.
 
  \item Implemented custom CUDA rasterization and defect-centered cropping operations bound to Python via \textbf{pybind11}, accelerating geometry-to-tensor rendering on the GPU and bypassing slow CPU-to-GPU memory copies.
 
  \item Engineered an asynchronous data loader using \textbf{NVIDIA DALI} to pipeline and overlap file prefetching and GPU tensor transfers with active deep learning inference via concurrent CUDA streams.
 
  \item Demonstrated that design-derived feature vectors performed on-par with legacy physical optical metrics, significantly improving the classification accuracy of production Gradient Boosted Decision Tree (GBDT) ensembles.
\end{itemize}
\vspace{-0.25em}

\textbf{Machine Learning Intern} \hfill May 2023 -- Aug 2023 \\
EMAG Technologies, Inc. \hfill \textit{Ann Arbor, MI}
\vspace{-0.5em}

\begin{itemize}\itemsep 0pt
  \item Built a closed-loop calibration application for active electronically scanned phased-array radar, optimizing 6-bit discrete phase-shifter and attenuator states on Anokiwave AWS-0103 beamforming ICs.
 
  \item Engineered the optimization loop around physical latency constraints, utilizing metaheuristic search algorithms (Genetic Algorithms / PSO) integrated with a software-state caching layer to prevent redundant SPI programming and SDR digitizing runs.
 
  \item Developed a CNN-based physical surrogate model that treated the 16$\times$16 aperture array as a 2D grid, learning mutual coupling and non-linear AM/PM distortion to predict optimal phase adjustments in software and bypass slow physical measurements.
 
  \item Achieved a \textbf{6.5$\times$ reduction} in post-calibration beam synthesis error compared to the baseline Sample Matrix Inversion (SMI) method under representative hardware-in-the-loop test configurations.
  \vspace{-0.25em}

\end{itemize}
\vspace{-0.5em}

\end{rSection}

\begin{rSection}{Projects} 
\vspace{-0.5em}

\textbf{Autonomous Fixed-Wing UAV with Object Classification} \hfill {\small C++, Python, OpenCV, ArduPilot, QGroundControl}

\vspace{-0.5em}
\begin{itemize}\itemsep -2pt
  \item Built a fully autonomous 1.6m fixed-wing UAV (MFD Crosswind Mini) for the AUVSI search-and-rescue competition; integrated Pixhawk Cube Orange, RTK GPS, dual-camera payload, and 900\,MHz telemetry into a CG-constrained airframe.
  \item Tuned ArduPilot's \textbf{EKF3 state estimation} and closed-loop control loops (\textbf{PID / TECS}) to achieve stable autonomous takeoff, waypoint-navigation, and failsafe maneuvers under variable wind conditions.
  \item Developed an Object Detection, Classification, and Localization (ODCL) pipeline utilizing a \textbf{YOLOv5} model to detect and geolocate targets in real-time from high-resolution Sony A6000 aerial imagery.
  \vspace{-0.25em}

\end{itemize}
\vspace{-0.5em}


\end{rSection}

\begin{rSection}{Selected Open-Source / Research}
\vspace{-0.5em}

\textbf{MetaRep (ICL $\leftrightarrow$ Kernel Ridge):} theory + experiments validating ridge-like ICL; reproducible harnesses and proofs. \\
\textbf{StableGLM (Rashomon sets for GLMs):} $\varepsilon$-Rashomon membership/certificates and set-level interpretability metrics. \\
\end{rSection}
\vspace{-0.5em}

\vspace{-0.5em}
\vspace{-0.5em}
\vspace{-0.5em}






%----------------------------------------------------------------------------------------
%	EXPERIENCE SECTION
%----------------------------------------------------------------------------------------



%----------------------------------------------------------------------------------------
%	PROJECTS SECTION
%----------------------------------------------------------------------------------------


%----------------------------------------------------------------------------------------
%	SKILLS SECTION
%----------------------------------------------------------------------------------------

\begin{rSection}{Skills}
\vspace{-0.2em}

\textbf{Programming Languages}: Python, C++, C, \text{C#}, TypeScript, Java, SQL, MATLAB, \LaTeX, Mojo
\vspace{-0.5em}

\textbf{HPC \& Acceleration}: CUDA, NVIDIA DALI, pybind11, Nsight Systems, gRPC / Protocol Buffers, GDSTK, OpenMP
\vspace{-0.5em}

\textbf{ML \& Infrastructure}: PyTorch, TensorFlow, Docker, Kubernetes, Docker Swarm, Postgres, Git, XGBoost / GBDTs
\vspace{-0.5em}

\textbf{Activities}: UM HuggingFace Admin, UM Hackers Quant Team, FPV Building/Racing, Archery 
\vspace{-1.5em}

\end{rSection}

\end{spacing}
\end{document}
